\documentclass[11pt]{article}
\usepackage[margin=1in]{geometry}
\usepackage{amsmath,amssymb,amsthm}
\usepackage{array,longtable,graphicx}
\usepackage[colorlinks=true,linkcolor=blue,citecolor=blue,urlcolor=blue]{hyperref}
\hypersetup{pdftitle={Unknown-frame recovery in sparse block-regular graphs},pdfauthor={Research manuscript}}
\newtheorem{theorem}{Theorem}[section]
\newtheorem{lemma}[theorem]{Lemma}
\newtheorem{proposition}[theorem]{Proposition}
\newtheorem{corollary}[theorem]{Corollary}
\newtheorem{definition}[theorem]{Definition}
\theoremstyle{remark}
\newtheorem{remark}[theorem]{Remark}
\newcommand{\one}{\mathbf 1}
\newcommand{\Prob}{\mathbb P}
\newcommand{\E}{\mathbb E}
\newcommand{\Fp}{\mathbb F_p}
\newcommand{\F}{\mathbb F}
\newcommand{\R}{\mathbb R}
\newcommand{\Q}{\mathbb Q}
\newcommand{\Z}{\mathbb Z}
\newcommand{\rank}{\operatorname{rank}}
\newcommand{\nullity}{\operatorname{nullity}}
\newcommand{\supp}{\operatorname{supp}}
\newcommand{\Var}{\operatorname{Var}}
\newcommand{\spec}{\operatorname{spec}}
\newcommand{\TV}{\operatorname{TV}}
\newcommand{\KL}{D_{\mathrm{KL}}}
\newcommand{\im}{\operatorname{im}}
\newcommand{\nnz}{\operatorname{nnz}}
\newcommand{\diag}{\operatorname{diag}}
\newcommand{\eps}{\varepsilon}
\renewcommand{\labelitemi}{$\bullet$}
\title{Unknown-frame recovery in sparse block-regular graphs\\Moments, certified algorithms, and sharp ambiguity}
\author{Research manuscript}
\date{November 15, 2026}
\begin{document}
\maketitle
\begin{abstract}
We recover an unknown bounded-size degree frame and its planted classes
in sparse block-regular graphs, from the graph and fixed order and degree
bounds alone. Independent uniform simple internal and cross factors may
have unequal class sizes, varying total degrees, and absent internal
edges. The hypothesis is a symmetric omission certificate with weights
in $[0,1/2)$ and unit physical row sums; zero weights are permitted.
The plant is, with high probability, the unique finest bounded-class
equitable partition. A bounded raw finite-field channel moment gives
complete partition listing in unconditional expected nearly quadratic
bit work and expected nearly linear space. Its proof uses heterogeneous
root/output spaces and selected cancellation-edge certificates to handle
inactive quotient channels. A certified complete-kernel algorithm
charges every input before averaging. We also prove disjoint old and
new spectra for all fixed connected simple regular base lifts, including
bipartite bases and irrational old modes. Exact two-class ambiguity and
hidden-prior Bayes error constants cover every $a\ge0$, $b\ge1$,
$a+b\ge3$. Complete bases give sharp lift ambiguity, while cycle lifts
have optimal success tending to zero. Separate prescribed-cross
spectral theorems yield rational recovery under their stated internal
degree hypotheses. Direct cover-recovery algorithms and classical
counting, mixing and linear-algebra inputs are credited; historical
priority is not certified.
\end{abstract}
\tableofcontents
\input{draft-v2/overview}
\input{draft-v2/weighted-unknown-frame}
\input{draft-v2/channel-moment-proof}
\input{draft-v2/scalar-frame-proof}
\input{draft-v2/full-kernel}
\input{draft-v2/vector-algorithm}
\input{draft-v2/lift-spectral}
\input{draft-v2/bipartite-spectral}
\input{draft-v2/lift-boundary}
\input{draft-v2/two-class-results}
\input{draft-v2/all-parameter-ambiguity}
\input{draft-v2/known-frame}
\input{draft-v2/fast-application}
\input{draft-v2/scope}
\appendix
\input{draft-v2/configuration-mixing}
\input{draft-v2/references}
\end{document}
